% =====================================================================
% Online Resource 1 (Electronic Supplementary Material) of the JIIS paper
% "OPENER: Open Partitioning Embedding for Named Entity Recognition".
% Upload ESM_1.pdf in "Supplementary material". The main text cites it as
% "Online Resource 1" (Tables S1-S2, Fig. S1).
% References to the main article (Table 5, Fig. 3, Section 4.7, ...) are read
% from ../main.aux with the xr package: compile the main article first.
% Compilation (from this folder): pdflatex ESM_1 && pdflatex ESM_1
% =====================================================================
\documentclass[10pt]{article}
\usepackage[a4paper, margin=2.5cm]{geometry}
\usepackage[utf8]{inputenc}
\usepackage[T1]{fontenc}
\usepackage{amsmath,amssymb}
\usepackage{graphicx}
\usepackage{xcolor}
\usepackage{colortbl}
\usepackage{booktabs}
\usepackage{threeparttable}
\usepackage{adjustbox}
\usepackage[labelfont=bf, labelsep=space, font=small]{caption}
\usepackage{xr}
\externaldocument[M-]{../main}

\definecolor{hlyellow}{rgb}{1,0.93,0.5}
\definecolor{hlred}{HTML}{FECACA}
\newcommand{\zstr}{OPENER-ZS\textsubscript{trans}}
\newcommand{\zsind}{OPENER-ZS\textsubscript{ind}}

% Supplementary numbering: Table S1, Fig. S1, ...
\renewcommand{\figurename}{Fig.}
\renewcommand{\thefigure}{S\arabic{figure}}
\renewcommand{\thetable}{S\arabic{table}}
\pagestyle{plain}

\begin{document}

% --- Required header (JIIS: title, journal, authors, affiliation and e-mail of the corresponding author)
\begin{center}
{\Large\bfseries Online Resource 1}\\[6pt]
{\large OPENER: Open Partitioning Embedding for Named Entity Recognition}\\[4pt]
\textit{Journal of Intelligent Information Systems}\\[6pt]
Thibault Garel, Faiza Belbachir, Assia Soukane\\
LyRIDS, ECE Engineering School, Paris, France\\
Corresponding author: Thibault Garel, \texttt{garel.thibault14@gmail.com}
\end{center}

\medskip
\noindent This Online Resource gathers supplementary tables and figures cited in the main article: the test-set bootstrap of the released model (Table~S1), the four-axis means of every system (Fig.~S1), and the zero-shot embedder comparison with Pile-NER (Table~S2).

% ------------------------------------------------------------------ Table S1
\begin{table}[h]
\centering
\caption{Test-set bootstrap robustness of the released (seed-$42$) model on three representative datasets}
\label{tab:bootstrap}
\small
\begin{adjustbox}{max width=\textwidth}
\begin{threeparttable}
\begin{tabular}{lccc}
\toprule
Model & CoNLL-2003 & CrossNER-Music & FabNER \\
\midrule
\zsind{}   & $66.1 \pm 1.4$ ~[$63.5$, $69.0$] & $55.5 \pm 0.9$ ~[$53.8$, $57.3$] & $36.2 \pm 0.7$ ~[$34.7$, $37.5$] \\
OPENER-Sup & $71.7 \pm 1.5$ ~[$68.9$, $74.6$] & $72.4 \pm 0.9$ ~[$70.6$, $74.2$] & $59.8 \pm 0.9$ ~[$57.9$, $61.6$] \\
\bottomrule
\end{tabular}
\begin{tablenotes}[flushleft]
\footnotesize
\item Typing-on-gold AMI ($\uparrow$) as bootstrap mean $\pm$ std [$95\%$ percentile CI], over $K{=}1000$ resamples of the gold mentions. The transductive refinement is excluded, as it adapts to the resampled test distribution. The bootstrap means match the full-test-set scores, and the stds stay at or below $1.5$ AMI, indicating stable scores.
\end{tablenotes}
\end{threeparttable}
\end{adjustbox}
\end{table}

% ------------------------------------------------------------------ Fig. S1
\begin{figure}[h]
\centering
\includegraphics{FigS1}
\caption{Per-model means on the four axes ($13$-set averages, Table~\ref{M-tab:efficiency} of the main article). The light dotted line separates the \emph{zero-shot} systems (left) from the \emph{supervised} OPENER-Sup (right). Latency, energy and CO\textsubscript{2} use a log scale. Labels: GL-S/M/L $=$ GLiNER sizes, OWN $=$ OWNER, ZS$_i$ $=$ \zsind, ZS$_t$ $=$ \zstr, Sup $=$ OPENER-Sup}
\label{fig:metrics-bars}
\end{figure}

% ------------------------------------------------------------------ Table S2
\begin{table}[h]
\centering
\caption{Zero-shot typing: effect of the contrastive embedder on \zsind (typing-on-gold, $13$-set means at the released seed ($42$))}
\label{tab:zs-embedder}
\small
\begin{adjustbox}{max width=\textwidth}
\begin{threeparttable}
\begin{tabular}{llccc}
\toprule
Embedder & Training data & AMI & F\textsubscript{1} & Acc \\
\midrule
hard-neg.\ mining (selected) & CoNLL + targets' train & \textbf{44.4} & 46.8 & 53.6 \\
Pile-NER (entity$\leftrightarrow$label) & Pile-NER only & 42.2 & \textbf{50.6} & \textbf{53.9} \\
\quad + label-name hard mining & Pile-NER only & 39.7 & 41.4 & 44.7 \\
\bottomrule
\end{tabular}
\begin{tablenotes}[flushleft]
\footnotesize
\item The nearest label-name centroid head is fixed, only the embedder changes. AMI, macro-$F_1$, accuracy ($\times100$, $\uparrow$). Best per column in bold. The selected hard-negative-mining embedder (target-tuned) is compared with Pile-NER variants (Pile-NER is the training corpus released with UniversalNER, cited in the main article) that are strictly domain-agnostic (never see target data). We report Pile-NER as a negative ablation (Section~\ref{M-sec:exp:ablations} of the main article).
\end{tablenotes}
\end{threeparttable}
\end{adjustbox}
\end{table}

\end{document}
